\section{Prime displacement terms and arbitrary fixed simple-even windows}
\label{sec:spectral-all-window-endpoint}

\begin{lemma}[AC-S42-kernel: bounded interior Fourier evaluations]
\label{AC-S42-kernel}
On a fixed interval \(I=(0,\ell)\), every \(\phi\in\mathcal X_I\)
has
\[
             \sup_N |P_N\phi(a)|<\infty
                         \qquad(0<a<\ell).               \tag{FW.1}
\]
The constant may depend on \(a,\ell,\phi\).  This is a bound
for the actual Fourier partial sums, not an assertion that point
evaluation is continuous on the logarithmic form domain.
\end{lemma}
\begin{proof}
Put \(k=2\pi/\ell\), and define
\[
 a(x)=\bigl(1+\tfrac12\log(1+k^2x^2)\bigr)^{-1},
 \qquad a_n=a(n),\qquad
 K_N(t)=\sum_{|n|\leq N}a_{|n|}e^{iknt}.                 \tag{FW.2}
\]
For large \(x\), direct differentiation gives
\[
 a(x)=O(1/\log x),\quad
 a'(x)=O(1/(x\log^2x)),\quad
 a''(x)=O(1/(x^2\log^2x)).                              \tag{FW.3}
\]
The function is smooth near zero and decreases to zero.
Let \(D_n\) and \(F_n=(n+1)^{-1}\sum_{j=0}^nD_j\) be
the periodic Dirichlet and Fejer kernels, and put
\(\Delta a_j=a_j-a_{j+1}\).
Two finite summations by parts give the exact identity
\[
 \begin{split}
 K_N={}&a_ND_N+\sum_{j=0}^{N-1}\Delta a_jD_j\\
     ={}&a_ND_N+N\Delta a_{N-1}F_{N-1}
             +\sum_{j=0}^{N-2}(j+1)\Delta^2a_jF_j.
 \end{split}                                             \tag{FW.4}
\]
Use normalized measure \(dt/\ell\) on the circle.
The kernels \(F_j\) are nonnegative and have integral one,
whereas \(\|D_N\|_1\leq C\log(N+2)\), by
\(|D_N(t)|\leq\min(2N+1,1/|\sin(\pi t/\ell)|)\).
Equation (FW.3) gives
\[
 \sup_N a_N\log(N+2)<\infty,\qquad
 \sup_N N|\Delta a_{N-1}|<\infty,\qquad
 \sum_{j\geq0}(j+1)|\Delta^2a_j|<\infty.                \tag{FW.5}
\]
For the last assertion use
\(\Delta^2a_j=\int_0^1\int_0^1a''(j+s+t)dsdt\)
and \(\int_0^\infty(x+1)|a''(x)|dx<\infty\).
Hence
\[
                          \sup_N\|K_N\|_1<\infty.        \tag{FW.6}
\]
Away from \(t=0\) modulo \(\ell\), the first identity in
(FW.4) also bounds \(K_N(t)\) uniformly in \(N\):
\(a_N+\sum_{j<N}\Delta a_j=a_0=1\), and the Dirichlet kernels
are uniformly bounded there.

By (FE.4), the function
\[
 h=(G_{\rm per}+1)\phi=G_I\phi-\mathcal K\phi+\phi       \tag{FW.7}
\]
belongs to \(L^2(I)\), and is bounded on every compact interior
subinterval.  Indeed \(G_I\phi\) and \(\phi\) are bounded,
and
\(\int_I K(x,y)dy=\int_{\mathbb R\setminus I}n(x-y)dy\)
is bounded when \(x\) stays away from the endpoints.
Fourier coefficients give
\[
                   P_N\phi(a)=\frac1\ell
                             \int_I K_N(a-y)h(y)dy.      \tag{FW.8}
\]
On a small neighborhood of \(a\), use the local boundedness
of \(h\) and (FW.6).  On its complement, use the uniform
off-diagonal kernel bound and \(h\in L^1(I)\).
Both bounds are independent of \(N\), proving (FW.1).
\end{proof}

\begin{theorem}[AC-S42: the finite normalization at every simple-even fixed window]
\label{AC-S42}
Fix any \(\lambda>1\) for which the actual closed Weil operator
has a simple even lowest eigenvalue.  Choose its real normalized
eigenvector \(u\) and the convergent real Galerkin phases \(u_N\).
Then there is a real \(\kappa_\lambda\ne0\) such that
\[
 \sqrt{\log N}\,u_N(\ell)\longrightarrow\kappa_\lambda,
 \qquad
 \frac{u_N/u_N(\ell)}{\sqrt{\log N}}
                   \longrightarrow\frac{u}{\kappa_\lambda}
                         \quad\hbox{in }L^2\hbox{ and form norm}.
                                                        \tag{FW.9}
\]
The exact coefficient is (FA.2).  The source determinants divided
by \(\sqrt{\log N}\) converge locally uniformly to the entire
function in (FA.4).  In particular the unrescaled determinants
diverge in magnitude at every \(z\) with \(\widehat u(z)\ne0\),
and cannot have a finite locally uniform entire limit.

No nonzero renewal trace of \(u\) is assumed.  The sole additional
spectral hypothesis is the simple-even lowest line stated above;
this theorem does not prove that hypothesis for all windows,
either unconditionally or from A.
\end{theorem}
\begin{proof}
Every part of the periodic comparison, translation gain, odd
compact-forcing construction and Ritz remainder in AC-S40 works
on an arbitrary fixed window with a simple even minimum.
In particular the barrier still has bounded \(A\)-image, because
all the additional prime translations are bounded on \(L^\infty\).
The odd resolvent has a positive gap from the lowest eigenvalue.
The only new term in the displacement calculation is the finite
prime sum.  The exact coefficient becomes
\[
 \begin{split}
 b_n={}&\frac1\pi\int_0^\ell(\mu(h)-2\cosh(h/2))
                                      \sin(2\pi nh/\ell)dh\\
       &+\frac1\pi\sum_{2\leq m\leq e^\ell}
            \frac{\Lambda(m)}{\sqrt m}
                                  \sin(2\pi n\log m/\ell).
 \end{split}                                             \tag{FW.10}
\]
This follows from the same source formula in
\texttt{basicexpli}, with the prime atoms in (EP.6).
The archimedean integral is still
\(\operatorname{sgn}(n)/4+O(1/|n|)\), so
\(\|\beta_N\|_2=O_\lambda(\sqrt N)\).

For an imaginary odd test \(\phi=iw\), its Fourier coefficients
\(\phi_n\) are real and odd.  For \(a\in(0,\ell)\),
\[
 \sum_{|n|\leq N}\phi_n\sin(2\pi na/\ell)
                           =-i\sqrt\ell\,P_N\phi(a).    \tag{FW.11}
\]
Each of these quantities is bounded by AC-S42-kernel.
If \(\log m=\ell\), its sine term is identically zero.
Thus all prime terms in \(\langle\beta_N,\phi\rangle\) are
\(O(1)\), and the exact leading Abelian coefficient remains
\[
 \frac{\langle\beta_N,\phi\rangle}{\sqrt{\log N}}
                      \longrightarrow
                      -\frac{\sqrt\ell}{\pi}\tau_+(w).  \tag{FW.12}
\]
The same fixed odd test with nonzero trace now proves
\(|u_N(\ell)|=O((\log N)^{-1/2})\) by (FE.25)--(FE.27).
Eventual finite endpoint nonvanishing follows, as before, from
the exact source displacement argument and even simplicity.
All the identities (FA.5)--(FA.9) are consequently valid on this
window.  Their all-test argument proves existence and nonvanishing
of \(\kappa_\lambda\), without using nonnegativity of \(u\).
Division by this nonzero limit proves (FW.9) and (FA.4).
Since \(u\ne0\), its compactly supported Fourier transform is
not identically zero.  At every point where it is nonzero,
(FA.4) gives the asserted divergence of the unrescaled determinant.
\end{proof}

\paragraph{Remaining fixed-window hypothesis.}
For a simple even minimum, prime translations no longer obstruct
the original finite endpoint asymptotic.  The unresolved
arbitrary-window assertion is the simple-even lowest-line property
itself.  AC-S43 identifies \(\kappa_\lambda\) exactly with the
nonzero renewal trace.  The large-window selection problem remains
a separate unresolved assertion.
